%% Generated by analysis/export_edges_tex.py. Do not edit by hand.
\begin{longtable}{@{}r l l l l p{0.40\textwidth}@{}}
\caption{The complete edge set of $K_G$. Tier codes: L lever, T treatable,
I indicator, M immutable. Grade is the evidence grade of the recommendation the
edge was extracted from; \emph{narr.} marks an edge drawn from section narrative
rather than from a numbered recommendation. 18 of the 26 edges carry a grade, over 21 distinct source--target pairs; the remaining 8 do not, and every edge entering the indicator tier is of that kind.}
\label{tab:edges}\\
\toprule
& source & target & grade & provenance & basis \\
\midrule
\endfirsthead
\multicolumn{6}{@{}l}{\footnotesize\itshape Table~\ref{tab:edges}, continued}\\
\toprule
& source & target & grade & provenance & basis \\
\midrule
\endhead
\midrule
\multicolumn{6}{r@{}}{\footnotesize\itshape continued on the next page}\\
\endfoot
\bottomrule
\endlastfoot
1 & Fruits\,(L) & HighChol\,(T) & C & SoC 2026, S10, Rec 10.15 & Lifestyle therapy is intensified for raised triglycerides or low HDL. \\
2 & PhysActivity\,(L) & HighChol\,(T) & C & SoC 2026, S10, Rec 10.15 & Lifestyle therapy is intensified for raised triglycerides or low HDL. \\
3 & Veggies\,(L) & HighChol\,(T) & C & SoC 2026, S10, Rec 10.15 & Lifestyle therapy is intensified for raised triglycerides or low HDL. \\
4 & BMI\,(L) & HighBP\,(T) & A & SoC 2026, S10, Rec 10.5 & Weight loss is the first lifestyle behaviour advised above 120/80 mmHg. \\
5 & Fruits\,(L) & HighBP\,(T) & A & SoC 2026, S10, Rec 10.5 & A DASH-style eating pattern is one of the advised lifestyle behaviours. \\
6 & HvyAlcoholConsump\,(L) & HighBP\,(T) & A & SoC 2026, S10, Rec 10.5 & Limiting or avoiding alcohol is one of the advised lifestyle behaviours. \\
7 & PhysActivity\,(L) & HighBP\,(T) & A & SoC 2026, S10, Rec 10.5 & Increased physical activity is one of the advised lifestyle behaviours. \\
8 & Smoker\,(L) & HighBP\,(T) & A & SoC 2026, S10, Rec 10.5 & Smoking cessation is one of the advised lifestyle behaviours. \\
9 & Veggies\,(L) & HighBP\,(T) & A & SoC 2026, S10, Rec 10.5 & A DASH-style eating pattern is one of the advised lifestyle behaviours. \\
10 & HighBP\,(T) & GenHlth\,(I) & narr. & SoC 2026, S10, narrative & Blood-pressure control is framed as reducing cardiovascular risk and improving outcomes. \\
11 & HighChol\,(T) & GenHlth\,(I) & narr. & SoC 2026, S10, narrative & Lipid control is framed as reducing cardiovascular risk and improving outcomes. \\
12 & Fruits\,(L) & BMI\,(L) & A & SoC 2026, S5, Rec 5.12 & Nutrition is one of the three components of the weight-management plan. \\
13 & PhysActivity\,(L) & BMI\,(L) & A & SoC 2026, S5, Rec 5.12 & A weight-management plan built on nutrition, physical activity and behavioural support targets 5--7\% weight loss. \\
14 & Veggies\,(L) & BMI\,(L) & A & SoC 2026, S5, Rec 5.12 & Nutrition is one of the three components of the weight-management plan. \\
15 & Fruits\,(L) & BMI\,(L) & B & SoC 2026, S5, Rec 5.14 & Eating patterns should emphasise whole fruits, among other principles. \\
16 & Veggies\,(L) & BMI\,(L) & B & SoC 2026, S5, Rec 5.14 & Eating patterns should emphasise nonstarchy vegetables, among other principles. \\
17 & HvyAlcoholConsump\,(L) & HighBP\,(T) & B & SoC 2026, S5, Rec 5.18 & Adults with or at risk of diabetes should not exceed recommended daily alcohol limits. \\
18 & PhysActivity\,(L) & GenHlth\,(I) & narr. & SoC 2026, S5, physical activity narrative & Physical activity is described as favourably affecting glycaemia, cardiovascular risk factors, weight, body composition and physical function. \\
19 & PhysActivity\,(L) & MentHlth\,(I) & narr. & SoC 2026, S5, physical activity narrative & Physical activity is described as favourably affecting psychosocial well-being, and as benefiting depressive symptoms. \\
20 & PhysActivity\,(L) & PhysHlth\,(I) & narr. & SoC 2026, S5, physical activity narrative & Physical activity is described as favourably affecting mobility and physical function. \\
21 & Smoker\,(L) & PhysHlth\,(I) & narr. & SoC 2026, S5, tobacco narrative & Smoking is described as causally linked to multiple health risks that affect morbidity in people with diabetes. \\
22 & BMI\,(L) & HighBP\,(T) & A & SoC 2026, S8, Rec 8.5 & Weight loss of 5--7\% improves glycaemia and intermediate cardiovascular risk factors. \\
23 & BMI\,(L) & HighChol\,(T) & A & SoC 2026, S8, Rec 8.5 & Same recommendation; dyslipidaemia is among the intermediate cardiovascular risk factors. \\
24 & PhysActivity\,(L) & BMI\,(L) & A & SoC 2026, S8, Rec 8.8a & Structured counselling centred on nutrition and physical activity to achieve a daily energy deficit. \\
25 & BMI\,(L) & GenHlth\,(I) & narr. & SoC 2026, S8, obesity narrative & Obesity treatment is framed as improving health outcomes broadly. \\
26 & BMI\,(L) & PhysHlth\,(I) & narr. & SoC 2026, S8, obesity narrative & Obesity treatment is framed as improving health outcomes broadly. \\
\end{longtable}
